\newif\ifglobal

%%%%%%%%%%%%%%%%%%%%%%%
%%% Compile to view %%%
%%%%%%%%%%%%%%%%%%%%%%%

%%%%%%%%%%%%%%%%%%%%%%%%%%%%%%%%%%%%%%%%%%%%%%%%%%%
%%% Readme:                                     %%%
%%% https://www.overleaf.com/read/bwdjvfjksvjy/ %%%
%%%%%%%%%%%%%%%%%%%%%%%%%%%%%%%%%%%%%%%%%%%%%%%%%%%

% >>> M?: marks a module setting number or important notice

% >>> M4: Set {./relative-path} to external files for this module <<<
\def \modulefiles {./31-RMT}
% To add an external file, use:
% (Replace '---' with actual filename)
% '\input{\modulefiles/tables/---.tex}' for tables
% '\includegraphics{\modulefiles/figures/---}' for figures
% '\subfile{\modulefiles/---}' for register files

% >>> M1: This module is in global project? <<<
% 'YES' - uncomment; 'NO' - comment out
\globaltrue

\ifglobal
    \documentclass[main\_\_EN.tex]{subfiles}

\else
    \documentclass[openany, 10pt]{book}

    % >>> M2: In ./chip-spec-settings/preamble-trm-module, also find
    %         settings for visibility of todo notes, watermark <<<
    \usepackage{preamble-trm-module}

    % >>> M3: Temporary labels in English,
    %         you can add your own labels there <<<
    \myexternaldocument{temp-labels-en}

\fi %\ifglobal

\setcounter{secnumdepth}{4}

% >>> M5: If this module needs any updates in
% './chip-spec-settings/preamble-trm-module.sty'
% (include additional package, etc.), contact your module PM
% or create the file './chip-spec-settings/changelog.tex' make the
% necessary updates yourself and document them in this file <<<

\raggedright


\begin{document}


% Sets language into which pre-defined bits of text are translated
\selectlanguage{english}

% Do not delete this block
\ifglobal
% Add the following front matter if
% this module is outside of global project
\else
    \InputIfFileExists{readme}{}{}
    {\let\clearpage\relax \listoftodos}
    {\let\clearpage\relax \listofdonetodos}
    \tableofcontents
    %\listoftables
    %\listoffigures
\fi


%%%%%%%%%%%%%%%%%%%%%%%%%
%%% TRM Module Begins %%%
%%%%%%%%%%%%%%%%%%%%%%%%%

\hypertarget{rmt}{}
\chapter{Remote Control Peripheral (RMT)}
\label{mod:rmt}


\section{Introduction}
The RMT (Remote Control) module is designed to send/receive infrared remote control signals that support for a variety of remote control protocols. The RMT module converts pulse codes stored in the module's built-in RAM into output signals, or converts input signals into pulse codes and stores them back in RAM. In addition, the RMT module optionally modulates its output signals with a carrier wave, or optionally filters its input signals.


\section{Functional Description}
\subsection{RMT Architecture}

The RMT module has four channels, numbered from zero to three. Each channel has the same functionality controlled by a dedicated set of registers and is able to transmit or receive data independently. By default, channel 0 transmits data from or receives data to block 0, channel 1 transmit data from or receive data to block 1, and so on, as shown in the figures below. Registers in each channel are indicated by \regindex{n} which is used as a placeholder for the channel number.

Each channel has a clock divider, counter, transmitter, and receiver. In each channel, only its transmitter or receiver can be active at the same time. The four channels share a 256 x 32-bit RAM.

\begin{figure}[H]
\includegraphics[width=1.0\textwidth]{\modulefiles/figures/31-RMT-architecture-20211014.eps}
\centering
\caption{RMT Architecture} \label{Figure:RMT2-1}
\end{figure}

\begin{figure}[H]
\includegraphics[width=0.8\textwidth]{\modulefiles/figures/31-RMT-channels-20211014.eps}
\centering
\caption{RMT Channels} \label{fig:rmt-channels}
\end{figure}



\subsection{RMT RAM}
\begin{figure}[htb! ]
\includegraphics[width=1.0\textwidth]{\modulefiles/figures/rmt_data_structure.eps}
\centering
\caption{Format of Pulse Code in RAM} \label{Figure:RMT2-2}
\end{figure}
The format of pulse code in RAM is shown in Figure \ref{Figure:RMT2-2}. Each pulse code contains a 16-bit entry with two fields, level and period. Level (0 or 1) indicates a high-/low-level value was received or is going to be sent, while "period" points out the clock cycles (see Figure \ref{Figure:RMT2-1} clk\_div) for which the level lasts. A zero period is interpreted as a transmission end-marker.

The RAM is divided into four 64 x 32-bit blocks. By default, each channel uses one block (block zero for channel zero, block one for channel one, and so on). Usually, only one block of 64 x 32-bit worth of data can be sent or received in channel \regindex{n}. If the data size is larger than this block size, users can configure the channel to enable wrap mode or to use more blocks by setting \hyperref[fielddesc:RMTREGMEMSIZECHN]{RMT\_MEM\_SIZE\_CH\regindex{n}}. Setting \hyperref[fielddesc:RMTREGMEMSIZECHN]{RMT\_MEM\_SIZE\_CH\regindex{n}} > 1 will prompt channel \regindex{n} to use the memory of subsequent channels, block (\regindex{n}) \verb+~+ block (\regindex{n} + \hyperref[fielddesc:RMTREGMEMSIZECHN]{RMT\_MEM\_SIZE\_CH\regindex{n}} -1). If so, the subsequent channels \regindex{n} +1 \verb+~+ \regindex{n} + \hyperref[fielddesc:RMTREGMEMSIZECHN]{RMT\_MEM\_SIZE\_CH\regindex{n}} - 1 cannot be used once their RAM blocks are occupied. Note that the RAM used by each channel is mapped from low address to high address. In such mode, channel 0 is able to use the RAM blocks for channels 1, 2 and 3 by setting \hyperref[fielddesc:RMTREGMEMSIZECHN]{RMT\_MEM\_SIZE\_CH\regindex{n}}, but channel 3 cannot use the blocks for channels 0, 1, or 2.

The RMT RAM can be accessed via APB bus, or read by the transmitter and written by the receiver. To protect a receiver from overwriting the blocks a transmitter is about to transmit, \hyperref[fielddesc:RMTREGMEMOWNERCHN]{RMT\_MEM\_OWNER\_CH\regindex{n}} can be configured to designate the block's owner, be it a transmitter or receiver. This way, if this ownership is violated, an \hyperref[fielddesc:RMTREGMEMOWNERERRCHN]{RMT\_MEM\_OWNER\_ERR\_CH\regindex{n}} flag will be generated.

When the RMT module is inactive, the RAM can be put into low-power mode by setting \hyperref[fielddesc:RMTREGMEMFORCEPD]{RMT\_MEM\_FORCE\_PD}.

\subsection{Clock}
The drive clock of a divider is generated by taking either the APB\_CLK or REF\_TICK according to the state of \hyperref[fielddesc:RMTREGREFALWAYSONCHN]{RMT\_REF\_ALWAYS\_ON\_CH\regindex{n}}. (For more information on clock sources, please see Chapter \hyperref[mod:resclk]{Reset and Clock}). Divider value is normally equal to the value of \hyperref[fielddesc:RMTREGDIVCNTCHN]{RMT\_DIV\_CNT\_CH\regindex{n}}, except value 0 that represents divider 256. The clock divider can be reset to zero by clearing \hyperref[fielddesc:RMTREGREFCNTRSTCHN]{RMT\_REF\_CNT\_RST\_CH\regindex{n}}. The clock generated from the divider can be used by the clock counter (see Figure \ref{Figure:RMT2-2}).

\subsection{Transmitter}
When \hyperref[fielddesc:RMTREGTXSTARTCHN]{RMT\_TX\_START\_CH\regindex{n}} is set to 1, the transmitter of channel \regindex{n} will start reading and sending pulse codes from the starting address of its RAM block. The codes are sent starting from low-address entry. The transmitter will stop the transmission, return to idle state and generate an \hyperref[int:RMTCHnTXENDINT]{RMT\_CH\regindex{n}\_TX\_END\_INT} interrupt, when an end-marker (a zero period) is encountered. Also, setting \hyperref[fielddesc:RMTREGTXSTOPCHN]{RMT\_TX\_STOP\_CH\regindex{n}} to 1 stops the transmission and immediately sets the transmitter back to idle. The output level of a transmitter in idle state is determined by the "level" field of the end-marker or by the content of \hyperref[fielddesc:RMTREGIDLEOUTLVCHN]{RMT\_IDLE\_OUT\_LV\_CH\regindex{n}}, depending on the configuration of \hyperref[fielddesc:RMTREGIDLEOUTENCHN]{RMT\_IDLE\_OUT\_EN\_CH\regindex{n}}.

To transmit more pulse codes than can be fitted in the channel's RAM, users can enable wrap mode by configuring \hyperref[fielddesc:RMTREGMEMTXWRAPEN]{RMT\_MEM\_TX\_WRAP\_EN}. In this mode, when the transmitter has reached the end-marker in the channel's memory, it will loop back to the first byte. For example, if \hyperref[fielddesc:RMTREGMEMSIZECHN]{RMT\_MEM\_SIZE\_CH\regindex{n}} is set to 1, the transmitter will start sending data from the address 64 * \regindex{n}, and then the data from higher RAM address. Once the transmitter finishes sending the data from (64 * (\regindex{n} +1) - 1), it will continue sending data from 64 * \regindex{n} till encounters an end-marker. Wrap mode is also applicable for \hyperref[fielddesc:RMTREGMEMSIZECHN]{RMT\_MEM\_SIZE\_CH\regindex{n}} > 1.

An \hyperref[int:RMTCHnTXTHREVENTINT]{RMT\_CH\regindex{n}\_TX\_THR\_EVENT\_INT} interrupt will be generated whenever the size of transmitted pulse codes is larger than or equal to the value set by \hyperref[fielddesc:RMTREGTXLIMCHN]{RMT\_TX\_LIM\_CH\regindex{n}}. In wrap mode, \hyperref[fielddesc:RMTREGTXLIMCHN]{RMT\_TX\_LIM\_CH\regindex{n}} can be set to a half or a fraction of the size of the channel's RAM block. When an \hyperref[int:RMTCHnTXTHREVENTINT]{RMT\_CH\regindex{n}\_TX\_THR\_EVENT\_INT} interrupt is detected, the already used RAM region should be updated by subsequent user defined events. Therefore, when the wrap mode happens the transmitter will seamlessly continue sending the new events.

The output of the transmitter can be modulated using a carrier wave by setting \hyperref[fielddesc:RMTREGCARRIERENCHN]{RMT\_CARRIER\_EN\_CH\regindex{n}}. The carrier waveform is configurable. In a carrier cycle, the high level lasts for (\hyperref[fielddesc:RMTREGCARRIERHIGHCHN]{RMT\_CARRIER\_HIGH\_CH\regindex{n}} +1) clock cycles of APB\_CLK or REF\_TICK, while the low level lasts for (\hyperref[fielddesc:RMTREGCARRIERLOWCHN]{RMT\_CARRIER\_LOW\_CH\regindex{n}} +1) clock cycles of APB\_CLK or REF\_TICK. When \hyperref[fielddesc:RMTREGCARRIEROUTLVCHN]{RMT\_CARRIER\_OUT\_LV\_CH\regindex{n}} is set to 1, carrier wave will be added on high-level output signals; while \hyperref[fielddesc:RMTREGCARRIEROUTLVCHN]{RMT\_CARRIER\_OUT\_LV\_CH\regindex{n}} is set to 0, carrier wave will be added on low-level output signals. Carrier wave can be added on output signals during modulation, or just added on valid pulse codes (the data stored in RAM), which can be set by configuring \hyperref[fielddesc:RMTREGCARRIEREFFENCHN]{RMT\_CARRIER\_EFF\_EN\_CH\regindex{n}}.

The continuous transmission of the transmitter can be enabled by setting \hyperref[fielddesc:RMTREGTXCONTIMODECHN]{RMT\_TX\_CONTI\_MODE\_CH\regindex{n}}. When this register is set, the transmitter will send the pulse codes from RAM in loops. If \hyperref[fielddesc:RMTREGTXLOOPCNTENCHN]{RMT\_TX\_LOOP\_CNT\_EN\_CH\regindex{n}} is set to 1, the transmitter will start counting loop times. Once the counting reaches the value of register \\
\hyperref[fielddesc:RMTREGTXLOOPNUMCHN]{RMT\_TX\_LOOP\_NUM\_CH\regindex{n}}, an \hyperref[int:RMTCHnTXLOOPINT]{RMT\_CH\regindex{n}\_TX\_LOOP\_INT} interrupt will be generated.

Setting \hyperref[fielddesc:RMTREGTXSIMEN]{RMT\_TX\_SIM\_EN} to 1 will enable multiple channels to start sending data simultaneously. \hyperref[fielddesc:RMTREGTXSIMCHN]{RMT\_TX\_SIM\_CH\regindex{n}} will choose which multiple channels are enabled to send data simultaneously.

\subsection{Receiver}
When \hyperref[fielddesc:RMTREGRXENCHN]{RMT\_RX\_EN\_CH\regindex{n}} is set to 1, the receiver in channel \regindex{n} becomes active, detecting signal levels and measuring clock cycles the signals lasts for. These data will be written in RAM in the form of pulse codes. Receiving ends, when the receiver detects no change in a signal level for a number of clock cycles more than the value set by \hyperref[fielddesc:RMTREGIDLETHRESCHN]{RMT\_IDLE\_THRES\_CH\regindex{n}}. The receiver will return to idle state and generate an \hyperref[int:RMTCHnRXENDINT]{RMT\_CH\regindex{n}\_RX\_END\_INT} interrupt.

The receiver has an input signal filter which can be enabled by configuring \hyperref[fielddesc:RMTREGRXFILTERENCHN]{RMT\_RX\_FILTER\_EN\_CH\regindex{n}}. The filter samples input signals continuously, and will detect the signals which remain unchanged for a continuous \hyperref[fielddesc:RMTREGRXFILTERTHRESCHN]{RMT\_RX\_FILTER\_THRES\_CH\regindex{n}} APB clock cycles as valid, otherwise, the signals will be detected as invalid. Only the valid signals can pass through this filter. The filter will remove pulses with a length of less than \hyperref[fielddesc:RMTREGRXFILTERTHRESCHN]{RMT\_RX\_FILTER\_THRES\_CH\regindex{n}} APB clock cycles.

\subsection{Interrupts}
\begin{itemize}
\item \label{int:RMTCHnERRINT}RMT\_CH\regindex{n}\_ERR\_INT: Triggered when channel \regindex{n} does not read or write data correctly. For example, if the transmitter still tries to read data from RAM when the RAM is empty, or the receiver still tries to write data into RAM when the RAM is full, this interrupt will be triggered.
\item \label{int:RMTCHnTXTHREVENTINT}RMT\_CH\regindex{n}\_TX\_THR\_EVENT\_INT: Triggered when the amount of data the transmitter has sent matches the value of \hyperref[fielddesc:RMTREGTXLIMCHN]{RMT\_TX\_LIM\_CH\regindex{n}}.
\item \label{int:RMTCHnTXENDINT}RMT\_CH\regindex{n}\_TX\_END\_INT: Triggered when the transmitter has finished transmitting signals.
\item \label{int:RMTCHnRXENDINT}RMT\_CH\regindex{n}\_RX\_END\_INT: Triggered when the receiver has finished receiving signals.
\item \label{int:RMTCHnTXLOOPINT}RMT\_CH\regindex{n}\_TX\_LOOP\_INT: Triggered when the loop counting reaches the value set by \hyperref[fielddesc:RMTREGTXLOOPNUMCHN]{RMT\_TX\_LOOP\_NUM\_CH\regindex{n}}.
\end{itemize}

\section{Base Address}\label{RMT_BASE_ADDRESS}
Users can access RMT with two base addresses, which can be seen in the following table. For more information about accessing peripherals from different buses please see Chapter \ref{mod:sysmem}: \textit{\nameref{mod:sysmem}}.

\begin{longtable}[c]{ | l | C{3cm} | }
\caption{RMT Base Address} \label{table:rmt_base_address}
\\ \hline
\rowcolor{lightgray}
Bus to Access Peripheral & Base Address \\
    \hline
PeriBUS1 & \hexprefix{}3F416000 \\
    \hline
PeriBUS2 & \hexprefix{}60016000 \\
    \hline
\end{longtable}

\hypertarget{rmt-reg-summ}{}
\section{Register Summary}

The addresses in the following table are relative to RMT base addresses provided in Section \ref{RMT_BASE_ADDRESS}.
\begin{longtable}{ | p{5.5cm} | p{7cm} | p{1.5cm} | p{1cm} | }
\hline\rowcolor{lightgray}
Name & Description &  Address & Access \\ \hline
\endhead
\multicolumn{4}{|l|}{\textbf{Configuration registers}} \\ \hline
\hyperref[regdesc:RMTCHNCONF0REG]{RMT\_CH0CONF0\_REG} & Channel 0 configuration register 0 & \hexprefix{}0010 & R/W \\ \hline
\hyperref[regdesc:RMTCHNCONF1REG]{RMT\_CH0CONF1\_REG} & Channel 0 configuration register 1 & \hexprefix{}0014 & varies \\ \hline
\hyperref[regdesc:RMTCHNCONF0REG]{RMT\_CH1CONF0\_REG} & Channel 1 configuration register 0 & \hexprefix{}0018 & R/W \\ \hline
\hyperref[regdesc:RMTCHNCONF1REG]{RMT\_CH1CONF1\_REG} & Channel 1 configuration register 1 & \hexprefix{}001C & varies \\ \hline
\hyperref[regdesc:RMTCHNCONF0REG]{RMT\_CH2CONF0\_REG} & Channel 2 configuration register 0 & \hexprefix{}0020 & R/W \\ \hline
\hyperref[regdesc:RMTCHNCONF1REG]{RMT\_CH2CONF1\_REG} & Channel 2 configuration register 1 & \hexprefix{}0024 & varies \\ \hline
\hyperref[regdesc:RMTCHNCONF0REG]{RMT\_CH3CONF0\_REG} & Channel 3 configuration register 0 & \hexprefix{}0028 & R/W \\ \hline
\hyperref[regdesc:RMTCHNCONF1REG]{RMT\_CH3CONF1\_REG} & Channel 3 configuration register 1 & \hexprefix{}002C & varies \\ \hline
\hyperref[regdesc:RMTAPBCONFREG]{RMT\_APB\_CONF\_REG} & RMT APB configuration register & \hexprefix{}0080 & R/W \\ \hline
\hyperref[regdesc:RMTREFCNTRSTREG]{RMT\_REF\_CNT\_RST\_REG} & RMT clock divider reset register & \hexprefix{}0088 & R/W \\ \hline
\hyperref[regdesc:RMTCHNRXCARRIERRMREG]{RMT\_CH0\_RX\_CARRIER\_RM\_REG} & Channel 0 carrier remove register & \hexprefix{}008C & R/W \\ \hline
\hyperref[regdesc:RMTCHNRXCARRIERRMREG]{RMT\_CH1\_RX\_CARRIER\_RM\_REG} & Channel 1 carrier remove register & \hexprefix{}0090 & R/W \\ \hline
\hyperref[regdesc:RMTCHNRXCARRIERRMREG]{RMT\_CH2\_RX\_CARRIER\_RM\_REG} & Channel 2 carrier remove register & \hexprefix{}0094 & R/W \\ \hline
\hyperref[regdesc:RMTCHNRXCARRIERRMREG]{RMT\_CH3\_RX\_CARRIER\_RM\_REG} & Channel 3 carrier remove register & \hexprefix{}0098 & R/W \\ \hline
\multicolumn{4}{|l|}{\textbf{Carrier wave duty cycle registers}} \\ \hline
\hyperref[regdesc:RMTCHNCARRIERDUTYREG]{RMT\_CH0CARRIER\_DUTY\_REG} & Channel 0 duty cycle configuration register & \hexprefix{}0060 & R/W \\ \hline
\hyperref[regdesc:RMTCHNCARRIERDUTYREG]{RMT\_CH1CARRIER\_DUTY\_REG} & Channel 1 duty cycle configuration register & \hexprefix{}0064 & R/W \\ \hline
\hyperref[regdesc:RMTCHNCARRIERDUTYREG]{RMT\_CH2CARRIER\_DUTY\_REG} & Channel 2 duty cycle configuration register & \hexprefix{}0068 & R/W \\ \hline
\hyperref[regdesc:RMTCHNCARRIERDUTYREG]{RMT\_CH3CARRIER\_DUTY\_REG} & Channel 3 duty cycle configuration register & \hexprefix{}006C & R/W \\ \hline
\multicolumn{4}{|l|}{\textbf{Tx event configuration registers}} \\ \hline
\hyperref[regdesc:RMTCHNTXLIMREG]{RMT\_CH0\_TX\_LIM\_REG} & Channel 0 Tx event configuration register & \hexprefix{}0070 & varies \\ \hline
\hyperref[regdesc:RMTCHNTXLIMREG]{RMT\_CH1\_TX\_LIM\_REG} & Channel 1 Tx event configuration register & \hexprefix{}0074 & varies \\ \hline
\hyperref[regdesc:RMTCHNTXLIMREG]{RMT\_CH2\_TX\_LIM\_REG} & Channel 2 Tx event configuration register & \hexprefix{}0078 & varies \\ \hline
\hyperref[regdesc:RMTCHNTXLIMREG]{RMT\_CH3\_TX\_LIM\_REG} & Channel 3 Tx event configuration register & \hexprefix{}007C & varies \\ \hline
\hyperref[regdesc:RMTTXSIMREG]{RMT\_TX\_SIM\_REG} & Enable RMT simultaneous transmission & \hexprefix{}0084 & R/W \\ \hline
\multicolumn{4}{|l|}{\textbf{Status registers}} \\ \hline
\hyperref[regdesc:RMTCHNSTATUSREG]{RMT\_CH0STATUS\_REG} & Channel 0 status register & \hexprefix{}0030 & RO \\ \hline
\hyperref[regdesc:RMTCHNSTATUSREG]{RMT\_CH1STATUS\_REG} & Channel 1 status register & \hexprefix{}0034 & RO \\ \hline
\hyperref[regdesc:RMTCHNSTATUSREG]{RMT\_CH2STATUS\_REG} & Channel 2 status register & \hexprefix{}0038 & RO \\ \hline
\hyperref[regdesc:RMTCHNSTATUSREG]{RMT\_CH3STATUS\_REG} & Channel 3 status register & \hexprefix{}003C & RO \\ \hline
\hyperref[regdesc:RMTCHNADDRREG]{RMT\_CH0ADDR\_REG} & Channel 0 address register & \hexprefix{}0040 & RO \\ \hline
\hyperref[regdesc:RMTCHNADDRREG]{RMT\_CH1ADDR\_REG} & Channel 1 address register & \hexprefix{}0044 & RO \\ \hline
\hyperref[regdesc:RMTCHNADDRREG]{RMT\_CH2ADDR\_REG} & Channel 2 address register & \hexprefix{}0048 & RO \\ \hline
\hyperref[regdesc:RMTCHNADDRREG]{RMT\_CH3ADDR\_REG} & Channel 3 address register & \hexprefix{}004C & RO \\ \hline
\multicolumn{4}{|l|}{\textbf{Version register}} \\ \hline
\hyperref[regdesc:RMTDATEREG]{RMT\_DATE\_REG} & Version control register & \hexprefix{}00FC & R/W \\ \hline
\multicolumn{4}{|l|}{\textbf{FIFO R/W registers}} \\ \hline
\hyperref[regdesc:RMTCHNDATAREG]{RMT\_CH0DATA\_REG} & Read and write data for channel 0 via APB FIFO & \hexprefix{}0000 & RO \\ \hline
\hyperref[regdesc:RMTCHNDATAREG]{RMT\_CH1DATA\_REG} & Read and write data for channel 1 via APB FIFO & \hexprefix{}0004 & RO \\ \hline
\hyperref[regdesc:RMTCHNDATAREG]{RMT\_CH2DATA\_REG} & Read and write data for channel 2 via APB FIFO & \hexprefix{}0008 & RO \\ \hline
\hyperref[regdesc:RMTCHNDATAREG]{RMT\_CH3DATA\_REG} & Read and write data for channel 3 via APB FIFO & \hexprefix{}000C & RO \\ \hline
\multicolumn{4}{|l|}{\textbf{Interrupt registers}} \\ \hline
\hyperref[regdesc:RMTINTRAWREG]{RMT\_INT\_RAW\_REG} & Raw interrupt status register & \hexprefix{}0050 & RO \\ \hline
\hyperref[regdesc:RMTINTSTREG]{RMT\_INT\_ST\_REG} & Masked interrupt status register & \hexprefix{}0054 & RO \\ \hline
\hyperref[regdesc:RMTINTENAREG]{RMT\_INT\_ENA\_REG} & Interrupt enable register & \hexprefix{}0058 & R/W \\ \hline
\hyperref[regdesc:RMTINTCLRREG]{RMT\_INT\_CLR\_REG} & Interrupt clear register & \hexprefix{}005C & WO \\ \hline
\end{longtable}


\section{Registers}
\begin{register}{H}{RMT\_CH\regindex{n}CONF0\_REG (\regindex{n}: 0-3)}{\hexprefix{}0010+8*\regindex{n}}\label{regdesc:RMTCHNCONF0REG}
\regfield{(reserved)}{2}{30}{00}%
\regfield{RMT\_CARRIER\_OUT\_LV\_CH\regindex{n}}{1}{29}{{1}}%
\regfield{RMT\_CARRIER\_EN\_CH\regindex{n}}{1}{28}{{1}}%
\regfield{RMT\_CARRIER\_EFF\_EN\_CH\regindex{n}}{1}{27}{{1}}%
\regfield{RMT\_MEM\_SIZE\_CH\regindex{n}}{3}{24}{{\hexprefix{}1}}%
\regfield{RMT\_IDLE\_THRES\_CH\regindex{n}}{16}{8}{{\hexprefix{}1000}}%
\regfield{RMT\_DIV\_CNT\_CH\regindex{n}}{8}{0}{{\hexprefix{}2}}%
\reglabel{Reset}\regnewline%
\begin{regdesc}\begin{reglist}
\label{fielddesc:RMTREGDIVCNTCHN}\item [RMT\_DIV\_CNT\_CH\regindex{n}] This field is used to configure clock divider for channel \regindex{n}. (R/W)
\label{fielddesc:RMTREGIDLETHRESCHN}\item [RMT\_IDLE\_THRES\_CH\regindex{n}] Receiving ends when no edge is detected on input signals for continuous clock cycles larger than this register value. (R/W)
\label{fielddesc:RMTREGMEMSIZECHN}\item [RMT\_MEM\_SIZE\_CH\regindex{n}] This field is used to configure the maximum blocks allocated to channel \regindex{n}. The valid range is from 1 \verb+~+ 4-\regindex{n}. (R/W)
\label{fielddesc:RMTREGCARRIEREFFENCHN}\item [RMT\_CARRIER\_EFF\_EN\_CH\regindex{n}] 1: Add carrier modulation on output signals only at data sending state for channel \regindex{n}. 0: Add carrier modulation on signals at all states for channel \regindex{n}. States here include idle state(ST\_IDLE), reading data from RAM (ST\_RD\_MEM), and sending data stored in RAM (ST\_SEND). Only valid when \hyperref[fielddesc:RMTREGCARRIERENCHN]{RMT\_CARRIER\_EN\_CH\regindex{n}} is set to 1. (R/W)
\label{fielddesc:RMTREGCARRIERENCHN}\item [RMT\_CARRIER\_EN\_CH\regindex{n}] This bit is used to enable carrier modulation for channel \regindex{n}. 1: Add carrier modulation on output signals. 0: No carrier modulation is added on output signals. (R/W)
\label{fielddesc:RMTREGCARRIEROUTLVCHN}\item [RMT\_CARRIER\_OUT\_LV\_CH\regindex{n}] This bit is used to configure the position of carrier wave for channel \regindex{n}.
1'h0: Add carrier wave on low-level output signals.
1'h1: Add carrier wave on high-level output signals. (R/W)
\end{reglist}\end{regdesc}
\end{register}


\begin{register}{H}{RMT\_CH\regindex{n}CONF1\_REG (\regindex{n}: 0-3)}{\hexprefix{}0014+8*\regindex{n}}\label{regdesc:RMTCHNCONF1REG}
\regfield{(reserved)}{11}{21}{00000000000}%
\regfield{RMT\_TX\_STOP\_CH\regindex{n}}{1}{20}{{0}}%
\regfield{RMT\_IDLE\_OUT\_EN\_CH\regindex{n}}{1}{19}{{0}}%
\regfield{RMT\_IDLE\_OUT\_LV\_CH\regindex{n}}{1}{18}{{0}}%
\regfield{RMT\_REF\_ALWAYS\_ON\_CH\regindex{n}}{1}{17}{{0}}%
\regfield{RMT\_CHK\_RX\_CARRIER\_EN\_CH\regindex{n}}{1}{16}{{0}}%
\regfield{RMT\_RX\_FILTER\_THRES\_CH\regindex{n}}{8}{8}{{\hexprefix{}f}}%
\regfield{RMT\_RX\_FILTER\_EN\_CH\regindex{n}}{1}{7}{{0}}%
\regfield{RMT\_TX\_CONTI\_MODE\_CH\regindex{n}}{1}{6}{{0}}%
\regfield{RMT\_MEM\_OWNER\_CH\regindex{n}}{1}{5}{{1}}%
\regfield{RMT\_APB\_MEM\_RST\_CH\regindex{n}}{1}{4}{{0}}%
\regfield{RMT\_MEM\_RD\_RST\_CH\regindex{n}}{1}{3}{{0}}%
\regfield{RMT\_MEM\_WR\_RST\_CH\regindex{n}}{1}{2}{{0}}%
\regfield{RMT\_RX\_EN\_CH\regindex{n}}{1}{1}{{0}}%
\regfield{RMT\_TX\_START\_CH\regindex{n}}{1}{0}{{0}}%
\reglabel{Reset}\regnewline%
\begin{regdesc}\begin{reglist}
\label{fielddesc:RMTREGTXSTARTCHN}\item [RMT\_TX\_START\_CH\regindex{n}] Set this bit to start sending data on channel \regindex{n}. (R/W)
\label{fielddesc:RMTREGRXENCHN}\item [RMT\_RX\_EN\_CH\regindex{n}] Set this bit to enable receiver to receive data on channel \regindex{n}. (R/W)
\label{fielddesc:RMTREGMEMWRRSTCHN}\item [RMT\_MEM\_WR\_RST\_CH\regindex{n}] Set this bit to reset RAM write address accessed by the receiver for channel \regindex{n}. (WO)
\label{fielddesc:RMTREGMEMRDRSTCHN}\item [RMT\_MEM\_RD\_RST\_CH\regindex{n}] Set this bit to reset RAM read address accessed by the transmitter for channel \regindex{n}. (WO)
\label{fielddesc:RMTREGMEMOWNERCHN}\item [RMT\_MEM\_OWNER\_CH\regindex{n}] This bit marks the ownership of channel \regindex{n}'s RAM block.
1'h1: Receiver is using the RAM.
1'h0: Transmitter is using the RAM. (R/W)
\label{fielddesc:RMTREGTXCONTIMODECHN}\item [RMT\_TX\_CONTI\_MODE\_CH\regindex{n}] Set this bit to restart transmission in continuous mode from the first data in channel \regindex{n}. (R/W)
\label{fielddesc:RMTREGRXFILTERENCHN}\item [RMT\_RX\_FILTER\_EN\_CH\regindex{n}] Set this bit to enable the receiver's filter for channel \regindex{n}. (R/W)
\label{fielddesc:RMTREGRXFILTERTHRESCHN}\item [RMT\_RX\_FILTER\_THRES\_CH\regindex{n}] Set this field to ignore the input pulse when its width is less than RMT\_RX\_FILTER\_THRES\_CH\regindex{n} APB clock cycles in receive mode. (R/W)
\label{fielddesc:RMTREGCHKRXCARRIERENCHN}\item [RMT\_CHK\_RX\_CARRIER\_EN\_CH\regindex{n}] Set this bit to enable memory loop read mode when carrier modulation is enabled for channel \regindex{n}. (R/W)
\label{fielddesc:RMTREGREFALWAYSONCHN}\item [RMT\_REF\_ALWAYS\_ON\_CH\regindex{n}] Set this bit to select a base clock for channel \regindex{n}.
1'h1: APB\_CLK; 1'h0: REF\_TICK (R/W)
\label{fielddesc:RMTREGIDLEOUTLVCHN}\item [RMT\_IDLE\_OUT\_LV\_CH\regindex{n}] This bit configures the level of output signals in channel \regindex{n} when the transmitter is in idle state. (R/W)
\label{fielddesc:RMTREGIDLEOUTENCHN}\item [RMT\_IDLE\_OUT\_EN\_CH\regindex{n}] This is the output enable bit for channel \regindex{n} in idle state. (R/W)
\label{fielddesc:RMTREGTXSTOPCHN}\item [RMT\_TX\_STOP\_CH\regindex{n}] Set this bit to stop the transmitter of channel \regindex{n} sending data out. (R/W)
\end{reglist}\end{regdesc}
\end{register}


\begin{register}{H}{RMT\_APB\_CONF\_REG}{\hexprefix{}0080}\label{regdesc:RMTAPBCONFREG}
\regfield{RMT\_CLK\_EN}{1}{31}{{0}}%
\regfield{(reserved)}{26}{5}{00000000000000000000000000}%
\regfield{RMT\_MEM\_FORCE\_PU}{1}{4}{{0}}%
\regfield{RMT\_MEM\_FORCE\_PD}{1}{3}{{0}}%
\regfield{RMT\_MEM\_CLK\_FORCE\_ON}{1}{2}{{1}}%
\regfield{RMT\_MEM\_TX\_WRAP\_EN}{1}{1}{{0}}%
\regfield{RMT\_APB\_FIFO\_MASK}{1}{0}{{0}}%
\reglabel{Reset}\regnewline%
\begin{regdesc}\begin{reglist}
\label{fielddesc:RMTREGAPBFIFOMASK}\item [RMT\_APB\_FIFO\_MASK] 1'h1: Access memory directly; 1'h0: Access memory via APB FIFO. (R/W)
\label{fielddesc:RMTREGMEMTXWRAPEN}\item [RMT\_MEM\_TX\_WRAP\_EN] Set this bit to enable wrap mode. (R/W)
\label{fielddesc:RMTREGMEMCLKFORCEON}\item [RMT\_MEM\_CLK\_FORCE\_ON] Set this bit to enable the clock for RAM when RMT module starts working; disable this clock when RMT stops working, to achieve low-power scheme. (R/W)
\label{fielddesc:RMTREGMEMFORCEPD}\item [RMT\_MEM\_FORCE\_PD] Set this bit to power down RMT memory. (R/W)
\label{fielddesc:RMTREGMEMFORCEPU}\item [RMT\_MEM\_FORCE\_PU] 1: Disable RAM's Light-sleep power down function. 0: power down RMT RAM when RMT is in Light-sleep mode. (R/W)
\label{fielddesc:RMTREGCLKEN}\item [RMT\_CLK\_EN] Clock gating enable bit for RMT registers to achieve low-power scheme. 1: Power up drive clock for RMT registers. 0: Power down drive clock for RMT registers. (R/W)
\end{reglist}\end{regdesc}
\end{register}


\begin{register}{H}{RMT\_REF\_CNT\_RST\_REG}{\hexprefix{}0088}\label{regdesc:RMTREFCNTRSTREG}
\regfield{(reserved)}{28}{4}{0000000000000000000000000000}%
\regfield{RMT\_REF\_CNT\_RST\_CH3}{1}{3}{{0}}%
\regfield{RMT\_REF\_CNT\_RST\_CH2}{1}{2}{{0}}%
\regfield{RMT\_REF\_CNT\_RST\_CH1}{1}{1}{{0}}%
\regfield{RMT\_REF\_CNT\_RST\_CH0}{1}{0}{{0}}%
\reglabel{Reset}\regnewline%
\begin{regdesc}\begin{reglist}
\label{fielddesc:RMTREGREFCNTRSTCHN}\item [RMT\_REF\_CNT\_RST\_CH\regindex{n}] This bit is used to reset the clock divider of channel \regindex{n}. (R/W)
\end{reglist}\end{regdesc}
\end{register}


\begin{register}{H}{RMT\_CH\regindex{n}\_RX\_CARRIER\_RM\_REG (\regindex{n}: 0-3)}{\hexprefix{}008C+4*\regindex{n}}\label{regdesc:RMTCHNRXCARRIERRMREG}
\regfield{RMT\_CARRIER\_HIGH\_THRES\_CH\regindex{n}}{16}{16}{{\hexprefix{}00}}%
\regfield{RMT\_CARRIER\_LOW\_THRES\_CH\regindex{n}}{16}{0}{{\hexprefix{}00}}%
\reglabel{Reset}\regnewline%
\begin{regdesc}\begin{reglist}
\label{fielddesc:RMTREGCARRIERLOWTHRESCHN}\item [RMT\_CARRIER\_LOW\_THRES\_CH\regindex{n}] The low level period in carrier modulation mode is (RMT\_CARRIER\_LOW\_THRES\_CH\regindex{n} + 1) clock cycles for channel \regindex{n}. (R/W)
\label{fielddesc:RMTREGCARRIERHIGHTHRESCHN}\item [RMT\_CARRIER\_HIGH\_THRES\_CH\regindex{n}] The high level period in carrier modulation mode is (RMT\_CARRIER\_HIGH\_THRES\_CH\regindex{n} + 1) clock cycles for channel \regindex{n}. (R/W)
\end{reglist}\end{regdesc}
\end{register}


\begin{register}{H}{RMT\_CH\regindex{n}CARRIER\_DUTY\_REG (\regindex{n}: 0-3)}{\hexprefix{}0060+4*\regindex{n}}\label{regdesc:RMTCHNCARRIERDUTYREG}
\regfield{RMT\_CARRIER\_HIGH\_CH\regindex{n}}{16}{16}{{\hexprefix{}40}}%
\regfield{RMT\_CARRIER\_LOW\_CH\regindex{n}}{16}{0}{{\hexprefix{}40}}%
\reglabel{Reset}\regnewline%
\begin{regdesc}\begin{reglist}
\label{fielddesc:RMTREGCARRIERLOWCHN}\item [RMT\_CARRIER\_LOW\_CH\regindex{n}] This field is used to configure the clock cycles of carrier wave at low level for channel \regindex{n}. (R/W)
\label{fielddesc:RMTREGCARRIERHIGHCHN}\item [RMT\_CARRIER\_HIGH\_CH\regindex{n}] This field is used to configure the clock cycles of carrier wave at high level for channel \regindex{n}. (R/W)
\end{reglist}\end{regdesc}
\end{register}


\begin{register}{H}{RMT\_CH\regindex{n}\_TX\_LIM\_REG (\regindex{n}: 0-3)}{\hexprefix{}0070+4*\regindex{n}}\label{regdesc:RMTCHNTXLIMREG}
\regfield{(reserved)}{11}{21}{00000000000}%
\regfield{RMT\_LOOP\_COUNT\_RESET\_CH\regindex{n}}{1}{20}{{0}}%
\regfield{RMT\_TX\_LOOP\_CNT\_EN\_CH\regindex{n}}{1}{19}{{0}}%
\regfield{RMT\_TX\_LOOP\_NUM\_CH\regindex{n}}{10}{9}{{\hexprefix{}0}}%
\regfield{RMT\_TX\_LIM\_CH\regindex{n}}{9}{0}{{\hexprefix{}80}}%
\reglabel{Reset}\regnewline%
\begin{regdesc}\begin{reglist}
\label{fielddesc:RMTREGTXLIMCHN}\item [RMT\_TX\_LIM\_CH\regindex{n}] This field is used to configure the maximum entries channel \regindex{n} can send out. When \hyperref[fielddesc:RMTREGMEMSIZECHN]{RMT\_MEM\_SIZE\_CH\regindex{n}} = 1, this field can be set to any value among 0 \verb+~+ 128 (64 * 32/16 = 128); when \hyperref[fielddesc:RMTREGMEMSIZECHN]{RMT\_MEM\_SIZE\_CH\regindex{n}} > 1, this field can be set to any value among (0 \verb+~+ 128) * \hyperref[fielddesc:RMTREGMEMSIZECHN]{RMT\_MEM\_SIZE\_CH\regindex{n}}. (R/W)
\label{fielddesc:RMTREGTXLOOPNUMCHN}\item [RMT\_TX\_LOOP\_NUM\_CH\regindex{n}] This field is used to configure the maximum loop times when continuous transmission mode is enabled. (R/W)
\label{fielddesc:RMTREGTXLOOPCNTENCHN}\item [RMT\_TX\_LOOP\_CNT\_EN\_CH\regindex{n}] This bit is used to enable loop counting. (R/W)
\label{fielddesc:RMTREGLOOPCOUNTRESETCHN}\item [RMT\_LOOP\_COUNT\_RESET\_CH\regindex{n}] This bit is used to reset loop counting when continuous transmission mode is valid. (WO)
\end{reglist}\end{regdesc}
\end{register}


\begin{register}{H}{RMT\_TX\_SIM\_REG}{\hexprefix{}0084}\label{regdesc:RMTTXSIMREG}
\regfield{(reserved)}{27}{5}{000000000000000000000000000}%
\regfield{RMT\_TX\_SIM\_EN}{1}{4}{{0}}%
\regfield{RMT\_TX\_SIM\_CH3}{1}{3}{{0}}%
\regfield{RMT\_TX\_SIM\_CH2}{1}{2}{{0}}%
\regfield{RMT\_TX\_SIM\_CH1}{1}{1}{{0}}%
\regfield{RMT\_TX\_SIM\_CH0}{1}{0}{{0}}%
\reglabel{Reset}\regnewline%
\begin{regdesc}\begin{reglist}
\label{fielddesc:RMTREGTXSIMCHN}\item [RMT\_TX\_SIM\_CH\regindex{n}] Set this bit to enable channel \regindex{n} to start sending data simultaneously with other enabled channels. (R/W)
\label{fielddesc:RMTREGTXSIMEN}\item [RMT\_TX\_SIM\_EN] This bit is used to enable multiple channels to start sending data simultaneously. (R/W)
\end{reglist}\end{regdesc}
\end{register}


\begin{register}{H}{RMT\_CH\regindex{n}STATUS\_REG (\regindex{n}: 0-3)}{\hexprefix{}0030+4*\regindex{n}}\label{regdesc:RMTCHNSTATUSREG}
\regfield{(reserved)}{4}{28}{0000}%
\regfield{RMT\_APB\_MEM\_RD\_ERR\_CH\regindex{n}}{1}{27}{{0}}%
\regfield{RMT\_APB\_MEM\_WR\_ERR\_CH\regindex{n}}{1}{26}{{0}}%
\regfield{RMT\_MEM\_EMPTY\_CH\regindex{n}}{1}{25}{{0}}%
\regfield{RMT\_MEM\_FULL\_CH\regindex{n}}{1}{24}{{0}}%
\regfield{RMT\_MEM\_OWNER\_ERR\_CH\regindex{n}}{1}{23}{{0}}%
\regfield{RMT\_STATE\_CH\regindex{n}}{3}{20}{{\hexprefix{}0}}%
\regfield{(reserved)}{1}{19}{0}%
\regfield{RMT\_MEM\_RADDR\_EX\_CH\regindex{n}}{9}{10}{{\hexprefix{}0}}%
\regfield{(reserved)}{1}{9}{0}%
\regfield{RMT\_MEM\_WADDR\_EX\_CH\regindex{n}}{9}{0}{{\hexprefix{}0}}%
\reglabel{Reset}\regnewline%
\begin{regdesc}\begin{reglist}
\label{fielddesc:RMTREGMEMWADDREXCHN}\item [RMT\_MEM\_WADDR\_EX\_CH\regindex{n}] This field records the memory address offset when receiver of channel \regindex{n} is using the RAM. (RO)
\label{fielddesc:RMTREGMEMRADDREXCHN}\item [RMT\_MEM\_RADDR\_EX\_CH\regindex{n}] This field records the memory address offset when transmitter of channel \regindex{n} is using the RAM. (RO)
\label{fielddesc:RMTREGSTATECHN}\item [RMT\_STATE\_CH\regindex{n}] This field records the FSM status of channel \regindex{n}. (RO)
\label{fielddesc:RMTREGMEMOWNERERRCHN}\item [RMT\_MEM\_OWNER\_ERR\_CH\regindex{n}] This status bit will be set when the ownership of memory block is violated. (RO)
\label{fielddesc:RMTREGMEMFULLCHN}\item [RMT\_MEM\_FULL\_CH\regindex{n}] This status bit will be set if the receiver receives more data than the memory allows. (RO)
\label{fielddesc:RMTREGMEMEMPTYCHN}\item [RMT\_MEM\_EMPTY\_CH\regindex{n}] This status bit will be set when the data to be sent is more than the memory allows and the wrap mode is disabled. (RO)
\label{fielddesc:RMTREGAPBMEMWRERRCHN}\item [RMT\_APB\_MEM\_WR\_ERR\_CH\regindex{n}] This status bit will be set if the offset address is out of memory size when channel \regindex{n} writes RAM via APB bus. (RO)
\label{fielddesc:RMTREGAPBMEMRDERRCHN}\item [RMT\_APB\_MEM\_RD\_ERR\_CH\regindex{n}] This status bit will be set if the offset address is out of memory size when channel \regindex{n} reads RAM via APB bus. (RO)
\end{reglist}\end{regdesc}
\end{register}


\begin{register}{H}{RMT\_CH\regindex{n}ADDR\_REG (\regindex{n}: 0-3)}{\hexprefix{}0040+4*\regindex{n}}\label{regdesc:RMTCHNADDRREG}
\regfield{(reserved)}{13}{19}{0000000000000}%
\regfield{RMT\_APB\_MEM\_RADDR\_CH\regindex{n}}{9}{10}{{\hexprefix{}0}}%
\regfield{(reserved)}{1}{9}{0}%
\regfield{RMT\_APB\_MEM\_WADDR\_CH\regindex{n}}{9}{0}{{\hexprefix{}0}}%
\reglabel{Reset}\regnewline%
\begin{regdesc}\begin{reglist}
\label{fielddesc:RMTREGAPBMEMWADDRCHN}\item [RMT\_APB\_MEM\_WADDR\_CH\regindex{n}] This field records the memory address offset when channel \regindex{n} writes RAM via APB bus. (RO)
\label{fielddesc:RMTREGAPBMEMRADDRCHN}\item [RMT\_APB\_MEM\_RADDR\_CH\regindex{n}] This field records the memory address offset when channel \regindex{n} reads RAM via APB bus. (RO)
\end{reglist}\end{regdesc}
\end{register}


\begin{register}{H}{RMT\_DATE\_REG}{\hexprefix{}00FC}\label{regdesc:RMTDATEREG}
\regfield{RMT\_DATE}{32}{0}{{\hexprefix{}19072601}}%
\reglabel{Reset}\regnewline%
\begin{regdesc}\begin{reglist}
\label{fielddesc:RMTREGDATE}\item [RMT\_DATE] Version control register. (R/W)
\end{reglist}\end{regdesc}
\end{register}


\begin{register}{H}{RMT\_CH\regindex{n}DATA\_REG (\regindex{n}: 0-3)}{\hexprefix{}0000+4*\regindex{n}}\label{regdesc:RMTCHNDATAREG}
\regfield{RMT\_CH\regindex{n}\_DATA}{32}{0}{{\hexprefix{}000000}}%
\reglabel{Reset}\regnewline%
\begin{regdesc}\begin{reglist}
\label{fielddesc:RMTCHNDATAREG}\item [RMT\_CH\regindex{n}\_DATA] This register is used to read and write data for channel \regindex{n} via APB FIFO. (RO)
\end{reglist}\end{regdesc}
\end{register}


\begin{register}{H}{RMT\_INT\_RAW\_REG}{\hexprefix{}0050}\label{regdesc:RMTINTRAWREG}
\regfield{(reserved)}{12}{20}{000000000000}%
\regfield{RMT\_CH3\_TX\_LOOP\_INT\_RAW}{1}{19}{{0}}%
\regfield{RMT\_CH2\_TX\_LOOP\_INT\_RAW}{1}{18}{{0}}%
\regfield{RMT\_CH1\_TX\_LOOP\_INT\_RAW}{1}{17}{{0}}%
\regfield{RMT\_CH0\_TX\_LOOP\_INT\_RAW}{1}{16}{{0}}%
\regfield{RMT\_CH3\_TX\_THR\_EVENT\_INT\_RAW}{1}{15}{{0}}%
\regfield{RMT\_CH2\_TX\_THR\_EVENT\_INT\_RAW}{1}{14}{{0}}%
\regfield{RMT\_CH1\_TX\_THR\_EVENT\_INT\_RAW}{1}{13}{{0}}%
\regfield{RMT\_CH0\_TX\_THR\_EVENT\_INT\_RAW}{1}{12}{{0}}%
\regfield{RMT\_CH3\_ERR\_INT\_RAW}{1}{11}{{0}}%
\regfield{RMT\_CH3\_RX\_END\_INT\_RAW}{1}{10}{{0}}%
\regfield{RMT\_CH3\_TX\_END\_INT\_RAW}{1}{9}{{0}}%
\regfield{RMT\_CH2\_ERR\_INT\_RAW}{1}{8}{{0}}%
\regfield{RMT\_CH2\_RX\_END\_INT\_RAW}{1}{7}{{0}}%
\regfield{RMT\_CH2\_TX\_END\_INT\_RAW}{1}{6}{{0}}%
\regfield{RMT\_CH1\_ERR\_INT\_RAW}{1}{5}{{0}}%
\regfield{RMT\_CH1\_RX\_END\_INT\_RAW}{1}{4}{{0}}%
\regfield{RMT\_CH1\_TX\_END\_INT\_RAW}{1}{3}{{0}}%
\regfield{RMT\_CH0\_ERR\_INT\_RAW}{1}{2}{{0}}%
\regfield{RMT\_CH0\_RX\_END\_INT\_RAW}{1}{1}{{0}}%
\regfield{RMT\_CH0\_TX\_END\_INT\_RAW}{1}{0}{{0}}%
\reglabel{Reset}\regnewline%
\begin{regdesc}\begin{reglist}
\label{fielddesc:RMTREGCHNTXENDINTRAW}\item [RMT\_CH\regindex{n}\_TX\_END\_INT\_RAW] The interrupt raw bit for channel \regindex{n}. Triggered when transmitting ends. (RO)
\label{fielddesc:RMTREGCHNRXENDINTRAW}\item [RMT\_CH\regindex{n}\_RX\_END\_INT\_RAW] The interrupt raw bit for channel \regindex{n}. Triggered when receiving ends. (RO)
\label{fielddesc:RMTREGCHNERRINTRAW}\item [RMT\_CH\regindex{n}\_ERR\_INT\_RAW] The interrupt raw bit for channel \regindex{n}. Triggered when error occurs. (RO)
\label{fielddesc:RMTREGCHNTXTHREVENTINTRAW}\item [RMT\_CH\regindex{n}\_TX\_THR\_EVENT\_INT\_RAW] The interrupt raw bit for channel \regindex{n}. Triggered when transmitter sends more data than configured value. (RO)
\label{fielddesc:RMTREGCHNTXLOOPINTRAW}\item [RMT\_CH\regindex{n}\_TX\_LOOP\_INT\_RAW] The interrupt raw bit for channel \regindex{n}. Triggered when loop counting reaches the configured threshold value. (RO)
\end{reglist}\end{regdesc}
\end{register}


\begin{register}{H}{RMT\_INT\_ST\_REG}{\hexprefix{}0054}\label{regdesc:RMTINTSTREG}
\regfield{(reserved)}{12}{20}{000000000000}%
\regfield{RMT\_CH3\_TX\_LOOP\_INT\_ST}{1}{19}{{0}}%
\regfield{RMT\_CH2\_TX\_LOOP\_INT\_ST}{1}{18}{{0}}%
\regfield{RMT\_CH1\_TX\_LOOP\_INT\_ST}{1}{17}{{0}}%
\regfield{RMT\_CH0\_TX\_LOOP\_INT\_ST}{1}{16}{{0}}%
\regfield{RMT\_CH3\_TX\_THR\_EVENT\_INT\_ST}{1}{15}{{0}}%
\regfield{RMT\_CH2\_TX\_THR\_EVENT\_INT\_ST}{1}{14}{{0}}%
\regfield{RMT\_CH1\_TX\_THR\_EVENT\_INT\_ST}{1}{13}{{0}}%
\regfield{RMT\_CH0\_TX\_THR\_EVENT\_INT\_ST}{1}{12}{{0}}%
\regfield{RMT\_CH3\_ERR\_INT\_ST}{1}{11}{{0}}%
\regfield{RMT\_CH3\_RX\_END\_INT\_ST}{1}{10}{{0}}%
\regfield{RMT\_CH3\_TX\_END\_INT\_ST}{1}{9}{{0}}%
\regfield{RMT\_CH2\_ERR\_INT\_ST}{1}{8}{{0}}%
\regfield{RMT\_CH2\_RX\_END\_INT\_ST}{1}{7}{{0}}%
\regfield{RMT\_CH2\_TX\_END\_INT\_ST}{1}{6}{{0}}%
\regfield{RMT\_CH1\_ERR\_INT\_ST}{1}{5}{{0}}%
\regfield{RMT\_CH1\_RX\_END\_INT\_ST}{1}{4}{{0}}%
\regfield{RMT\_CH1\_TX\_END\_INT\_ST}{1}{3}{{0}}%
\regfield{RMT\_CH0\_ERR\_INT\_ST}{1}{2}{{0}}%
\regfield{RMT\_CH0\_RX\_END\_INT\_ST}{1}{1}{{0}}%
\regfield{RMT\_CH0\_TX\_END\_INT\_ST}{1}{0}{{0}}%
\reglabel{Reset}\regnewline%
\begin{regdesc}\begin{reglist}
\label{fielddesc:RMTREGCHNTXENDINTST}\item [RMT\_CH\regindex{n}\_TX\_END\_INT\_ST] The masked interrupt status bit for \hyperref[int:RMTCHnTXENDINT]{RMT\_CH\regindex{n}\_TX\_END\_INT}. (RO)
\label{fielddesc:RMTREGCHNRXENDINTST}\item [RMT\_CH\regindex{n}\_RX\_END\_INT\_ST] The masked interrupt status bit for \hyperref[int:RMTCHnRXENDINT]{RMT\_CH\regindex{n}\_RX\_END\_INT}. (RO)
\label{fielddesc:RMTREGCHNERRINTST}\item [RMT\_CH\regindex{n}\_ERR\_INT\_ST] The masked interrupt status bit for \hyperref[int:RMTCHnERRINT]{RMT\_CH\regindex{n}\_ERR\_INT}. (RO)
\label{fielddesc:RMTREGCHNTXTHREVENTINTST}\item [RMT\_CH\regindex{n}\_TX\_THR\_EVENT\_INT\_ST] The masked interrupt status bit for \hyperref[int:RMTCHnTXTHREVENTINT]{RMT\_CH\regindex{n}\_TX\_THR\_EVENT\_INT}. (RO)
\label{fielddesc:RMTREGCHNTXLOOPINTST}\item [RMT\_CH\regindex{n}\_TX\_LOOP\_INT\_ST] The masked interrupt status bit for \hyperref[int:RMTCHnTXLOOPINT]{RMT\_CH\regindex{n}\_TX\_LOOP\_INT}. (RO)
\end{reglist}\end{regdesc}
\end{register}


\begin{register}{H}{RMT\_INT\_ENA\_REG}{\hexprefix{}0058}\label{regdesc:RMTINTENAREG}
\regfield{(reserved)}{12}{20}{000000000000}%
\regfield{RMT\_CH3\_TX\_LOOP\_INT\_ENA}{1}{19}{{0}}%
\regfield{RMT\_CH2\_TX\_LOOP\_INT\_ENA}{1}{18}{{0}}%
\regfield{RMT\_CH1\_TX\_LOOP\_INT\_ENA}{1}{17}{{0}}%
\regfield{RMT\_CH0\_TX\_LOOP\_INT\_ENA}{1}{16}{{0}}%
\regfield{RMT\_CH3\_TX\_THR\_EVENT\_INT\_ENA}{1}{15}{{0}}%
\regfield{RMT\_CH2\_TX\_THR\_EVENT\_INT\_ENA}{1}{14}{{0}}%
\regfield{RMT\_CH1\_TX\_THR\_EVENT\_INT\_ENA}{1}{13}{{0}}%
\regfield{RMT\_CH0\_TX\_THR\_EVENT\_INT\_ENA}{1}{12}{{0}}%
\regfield{RMT\_CH3\_ERR\_INT\_ENA}{1}{11}{{0}}%
\regfield{RMT\_CH3\_RX\_END\_INT\_ENA}{1}{10}{{0}}%
\regfield{RMT\_CH3\_TX\_END\_INT\_ENA}{1}{9}{{0}}%
\regfield{RMT\_CH2\_ERR\_INT\_ENA}{1}{8}{{0}}%
\regfield{RMT\_CH2\_RX\_END\_INT\_ENA}{1}{7}{{0}}%
\regfield{RMT\_CH2\_TX\_END\_INT\_ENA}{1}{6}{{0}}%
\regfield{RMT\_CH1\_ERR\_INT\_ENA}{1}{5}{{0}}%
\regfield{RMT\_CH1\_RX\_END\_INT\_ENA}{1}{4}{{0}}%
\regfield{RMT\_CH1\_TX\_END\_INT\_ENA}{1}{3}{{0}}%
\regfield{RMT\_CH0\_ERR\_INT\_ENA}{1}{2}{{0}}%
\regfield{RMT\_CH0\_RX\_END\_INT\_ENA}{1}{1}{{0}}%
\regfield{RMT\_CH0\_TX\_END\_INT\_ENA}{1}{0}{{0}}%
\reglabel{Reset}\regnewline%
\begin{regdesc}\begin{reglist}
\label{fielddesc:RMTREGCHNTXENDINTENA}\item [RMT\_CH\regindex{n}\_TX\_END\_INT\_ENA] Interrupt enable bit for \hyperref[int:RMTCHnTXENDINT]{RMT\_CH\regindex{n}\_TX\_END\_INT}. (R/W)
\label{fielddesc:RMTREGCHNRXENDINTENA}\item [RMT\_CH\regindex{n}\_RX\_END\_INT\_ENA] Interrupt enable bit for \hyperref[int:RMTCHnRXENDINT]{RMT\_CH\regindex{n}\_RX\_END\_INT}. (R/W)
\label{fielddesc:RMTREGCHNERRINTENA}\item [RMT\_CH\regindex{n}\_ERR\_INT\_ENA] Interrupt enable bit for \hyperref[int:RMTCHnERRINT]{RMT\_CH\regindex{n}\_ERR\_INT}. (R/W)
\label{fielddesc:RMTREGCHNTXTHREVENTINTENA}\item [RMT\_CH\regindex{n}\_TX\_THR\_EVENT\_INT\_ENA] Interrupt enable bit for \hyperref[int:RMTCHnTXTHREVENTINT]{RMT\_CH\regindex{n}\_TX\_THR\_EVENT\_INT}. (R/W)
\label{fielddesc:RMTREGCHNTXLOOPINTENA}\item [RMT\_CH\regindex{n}\_TX\_LOOP\_INT\_ENA] Interrupt enable bit for \hyperref[int:RMTCHnTXLOOPINT]{RMT\_CH\regindex{n}\_TX\_LOOP\_INT}. (R/W)
\end{reglist}\end{regdesc}
\end{register}


\begin{register}{H}{RMT\_INT\_CLR\_REG}{\hexprefix{}005C}\label{regdesc:RMTINTCLRREG}
\regfield{(reserved)}{12}{20}{000000000000}%
\regfield{RMT\_CH3\_TX\_LOOP\_INT\_CLR}{1}{19}{{0}}%
\regfield{RMT\_CH2\_TX\_LOOP\_INT\_CLR}{1}{18}{{0}}%
\regfield{RMT\_CH1\_TX\_LOOP\_INT\_CLR}{1}{17}{{0}}%
\regfield{RMT\_CH0\_TX\_LOOP\_INT\_CLR}{1}{16}{{0}}%
\regfield{RMT\_CH3\_TX\_THR\_EVENT\_INT\_CLR}{1}{15}{{0}}%
\regfield{RMT\_CH2\_TX\_THR\_EVENT\_INT\_CLR}{1}{14}{{0}}%
\regfield{RMT\_CH1\_TX\_THR\_EVENT\_INT\_CLR}{1}{13}{{0}}%
\regfield{RMT\_CH0\_TX\_THR\_EVENT\_INT\_CLR}{1}{12}{{0}}%
\regfield{RMT\_CH3\_ERR\_INT\_CLR}{1}{11}{{0}}%
\regfield{RMT\_CH3\_RX\_END\_INT\_CLR}{1}{10}{{0}}%
\regfield{RMT\_CH3\_TX\_END\_INT\_CLR}{1}{9}{{0}}%
\regfield{RMT\_CH2\_ERR\_INT\_CLR}{1}{8}{{0}}%
\regfield{RMT\_CH2\_RX\_END\_INT\_CLR}{1}{7}{{0}}%
\regfield{RMT\_CH2\_TX\_END\_INT\_CLR}{1}{6}{{0}}%
\regfield{RMT\_CH1\_ERR\_INT\_CLR}{1}{5}{{0}}%
\regfield{RMT\_CH1\_RX\_END\_INT\_CLR}{1}{4}{{0}}%
\regfield{RMT\_CH1\_TX\_END\_INT\_CLR}{1}{3}{{0}}%
\regfield{RMT\_CH0\_ERR\_INT\_CLR}{1}{2}{{0}}%
\regfield{RMT\_CH0\_RX\_END\_INT\_CLR}{1}{1}{{0}}%
\regfield{RMT\_CH0\_TX\_END\_INT\_CLR}{1}{0}{{0}}%
\reglabel{Reset}\regnewline%
\begin{regdesc}\begin{reglist}
\label{fielddesc:RMTREGCHNTXENDINTCLR}\item [RMT\_CH\regindex{n}\_TX\_END\_INT\_CLR] Set this bit to clear \hyperref[int:RMTCHnTXENDINT]{RMT\_CH\regindex{n}\_TX\_END\_INT} interrupt. (WO)
\label{fielddesc:RMTREGCHNRXENDINTCLR}\item [RMT\_CH\regindex{n}\_RX\_END\_INT\_CLR] Set this bit to clear \hyperref[int:RMTCHnRXENDINT]{RMT\_CH\regindex{n}\_RX\_END\_INT} interrupt. (WO)
\label{fielddesc:RMTREGCHNERRINTCLR}\item [RMT\_CH\regindex{n}\_ERR\_INT\_CLR] Set this bit to clear \hyperref[int:RMTCHnERRINT]{RMT\_CH\regindex{n}\_ERR\_INT} interrupt. (WO)
\label{fielddesc:RMTREGCHNTXTHREVENTINTCLR}\item [RMT\_CH\regindex{n}\_TX\_THR\_EVENT\_INT\_CLR] Set this bit to clear \hyperref[int:RMTCHnTXTHREVENTINT]{RMT\_CH\regindex{n}\_TX\_THR\_EVENT\_INT} interrupt. (WO)
\label{fielddesc:RMTREGCHNTXLOOPINTCLR}\item [RMT\_CH\regindex{n}\_TX\_LOOP\_INT\_CLR] Set this bit to clear \hyperref[int:RMTCHnTXLOOPINT]{RMT\_CH\regindex{n}\_TX\_LOOP\_INT} interrupt. (WO)
\end{reglist}\end{regdesc}
\end{register}


\end{document}
